\documentclass[10pt,a4paper]{amsart}
\usepackage[margin=19mm]{geometry}
\usepackage{amsmath,amssymb,mathtools}
\usepackage[scr=rsfs,cal=euler]{mathalfa}
\usepackage{xfrac}
\usepackage[unicode,hidelinks,bookmarks=false]{hyperref}
\newcommand{\ie}{i.e.\ }
\newcommand{\cf}{cf.\ }
\newcommand{\resp}{respectively\ }
\newcommand{\Subsection}{\subsection}
\providecommand{\nobreakdash}{\nobreak\hbox{-}}
\setlength{\emergencystretch}{2em}
\multlinegap0pt
\usepackage{amssymb}
\usepackage{bm}
\usepackage{float}
\usepackage{mathtools}
\newtheorem{lem}{Lemma}
\title{Cubic Polynomials and Sums of Two Squares}
\author{Siddharth Iyer}
\date{Source: arXiv:2510.25492v3 (CC BY 4.0)}
\begin{document}
\maketitle

\noindent\emph{Source and licence.} Cubic Polynomials and Sums of Two Squares. Version arXiv:2510.25492v3 (CC BY 4.0), distributed by arXiv under the Creative Commons Attribution 4.0 International licence, http://creativecommons.org/licenses/by/4.0/, as recorded in the arXiv licence metadata for this exact version and shown on the abstract page https://arxiv.org/abs/2510.25492v3. Full licence notice: \texttt{licenses/arxiv-2510.25492-CC-BY-4.0.txt}.\par\smallskip

\noindent\textit{Editorial scope.} Counted unit: the article's lem betaalphaauxlemmasum2squares (source0.tex lines 781--791). The article's complete original proof of it is delivered with it (source0.tex lines 792--871, one \texttt{proof} environment of 80 lines). No mathematics is written, repaired or re-derived: the statement and the proof are the article's own lines, in the article's own notation, and no retained line is substituted or re-flowed.  Retained uncounted as the source-local context the proof needs: lines 689--692.  Nothing was omitted from any retained span, so the package carries no omission record.  No retained line contains a citation, so the wrapper carries no bibliography and no bibliography entry.  No retained line contains a figure environment, an image-inclusion command or drawing code, so no image is delivered and the package declares no asset dependency.  Editorial formatting only, in addition: an AMS \texttt{amsart} preamble that restates the article's own declarations and macros used by the retained lines, namely the environments \texttt{lem} on separate counters, together with the standard packages those lines need.  The retained environments print in this document as Lemma 1 in order of appearance, whereas the article prints the same items with the numbering of its own full document.  The complete native source of the article as distributed in the arXiv e-print for this exact version is bundled unchanged as \texttt{sources/188041/source0.tex}.
\begin{lem}
\label{SsubsetRtheta1inftyinfty}
We have $\mathbf{S} \subset \mathcal{R}_{\theta,1}(\infty,\infty)$.
\end{lem}

\begin{lem}
\label{betaalphaauxlemmasum2squares}
If $c>0$, there exists a positive constant $F(c)$ with $\lim_{c\rightarrow 0^{+}}F(c) = 0$ such that if $X$ is sufficiently large ($X \geq N_{c}$) and 
\begin{align*}
1\leq \beta \leq \alpha \leq c\cdot X^{1/6},
\end{align*}
we have  
\begin{align*}
\mathbf{S}_{c}(\alpha,\beta,X) \subset \mathcal{R}_{\theta,1}(F(c)\cdot X^{1/2},\infty).
\end{align*}
\end{lem}

\begin{proof}
If
\begin{align*}
((u_{0}+u_{1}\theta+u_{2}\theta^2), (v_{0}+v_{1}\theta+v_{2}\theta^2)) \in \mathbf{S}_{c}(\alpha,\beta,X),
\end{align*}
by Lemma \ref{SsubsetRtheta1inftyinfty} we know that 
\begin{align}
\label{Halfdoneu0u1u2v0v1v2proof}
((u_{0}+u_{1}\theta+u_{2}\theta^2), (v_{0}+v_{1}\theta+v_{2}\theta^2)) \in \mathcal{R}_{\theta,1}(\infty,\infty).
\end{align}
It remains to show that $|u_{i}|,|v_{i}| \leq F(c)\cdot X^{1/2}$ for some function $F$ which satisfies $\lim_{c\rightarrow 0^{+}}F(c) = 0$. Note that $I(u_{1},\beta,v_{1},\alpha) = 1$ or 
\begin{align*}
u_{1}\alpha-v_{1}\beta = 1,
\end{align*}
implying that
\begin{align*}
|u_{1}| = \left|\frac{1+v_{1}\beta}{\alpha}\right| \leq 2 |v_{1}|.
\end{align*}
Using
\begin{align}
\label{v1inequalitysum2squares}
|v_{1}|\leq c\cdot X^{1/6},
\end{align}
we obtain
\begin{align}
\label{u1inequalitysum2squares}
|u_{1}| \leq 2c \cdot X^{1/6}.
\end{align}
On the set $\mathbf{S}_{c}(\alpha,\beta,X),$ we thus have
\begin{align*}
&h_{1}(u_{1},u_{2}, v_{1},v_{2}) = h_{1}(u_{1},\beta, v_{1},\alpha)\\
&= -1+2a_{1}u_{1}\beta-\beta^2(a_{1}a_{2}-a_{0})+2a_{1}v_{1}\alpha-\alpha^2(a_{1}a_{2}-a_{0}),
\end{align*}
giving us the estimate
\begin{align*}
\left|h_{1}(u_{1},u_{2}, v_{1},v_{2})\right| \leq 1+2|a_{1}|c^2 X^{1/3} + 2c^2|(a_{1}a_{2}-a_{0})|X^{1/3}+2c^2|a_{1}|X^{1/3},
\end{align*}
so that if $X$ is sufficiently large we have 
\begin{align}
\label{h1u1u2v1v2inequalitysum2squares}
\left|h_{1}(u_{1},u_{2}, v_{1},v_{2})\right| &\leq \left(4|a_{1}|c^2 + 4c^2|(a_{1}a_{2}-a_{0})|+4c^2|a_{1}|\right)X^{1/3}\notag\\
&= H_{1}(c)\cdot X^{1/3},
\end{align}
where $H_{1}(c)>0$ is some constant that approaches $0$ as $c$ approaches $0$.
Similarly we obtain
\begin{align*}
&h_{2}(u_{1},u_{2}, v_{1},v_{2}) = h_{2}(u_{1},\beta, v_{1},\alpha)\\
&= -u_{1}^2+2a_{2}u_{1}\beta-\beta^2(a_{2}^2-a_{1}) -v_{1}^2+2a_{2}\alpha v_{2}-\alpha^2(a_{2}^2-a_{1}),
\end{align*}
so that with \eqref{v1inequalitysum2squares} and \eqref{u1inequalitysum2squares} we have
\begin{align}
\label{h2u1u2v1v2inequalitysum2squares}
|h_{2}(u_{1},u_{2}, v_{1},v_{2})| \leq H_{2}(c)\cdot X^{1/3},
\end{align}
where $H_{2}(c)>0$ is some constant that approaches $0$ as $c$ approaches $0$. Note that 
\begin{align*}
u_{0} &= U_{0}(u_{1},u_{2},v_{1},v_{2}) \\
&= v_{2} h_{1}(u_{1},u_{2},v_{1},v_{2})/2 - v_{1} h_{2}(u_{1},u_{2},v_{1},v_{2})/2.
\end{align*}
Hence
\begin{align*}
|u_{0}| \leq |v_{2} h_{1}(u_{1},u_{2},v_{1},v_{2})| +|v_{1} h_{2}(u_{1},u_{2},v_{1},v_{2})|,
\end{align*}
Using inequalities \eqref{v1inequalitysum2squares}, \eqref{u1inequalitysum2squares}, \eqref{h1u1u2v1v2inequalitysum2squares}, \eqref{h2u1u2v1v2inequalitysum2squares} and substitutions $v_{2} = \alpha$ and $u_{2} = \beta$, we obtain 
\begin{align}
\label{u0inequalitysum2squares}
|u_{0}| \leq \Theta(c)\cdot X^{1/2},
\end{align}
where $\Theta(c)>0$ is some constant that approaches $0$ as $c$ approaches $0$. Similarly we obtain
\begin{align*}
v_{0} &= V_{0}(u_{1},u_{2},v_{1},v_{2})\\
&= -u_{2} h_{1}(u_{1},u_{2},v_{1},v_{2})/2 +u_{1}h_{2}(u_{1},u_{2},v_{1},v_{2})/2.
\end{align*}
Using inequalities \eqref{v1inequalitysum2squares}, \eqref{u1inequalitysum2squares}, \eqref{h1u1u2v1v2inequalitysum2squares}, \eqref{h2u1u2v1v2inequalitysum2squares} and substitutions $v_{2} = \alpha$ and $u_{2} = \beta$ we obtain
\begin{align}
\label{v0inequalitysum2squares}
|v_{0}| \leq \rho(c)\cdot X^{1/2},
\end{align}
where $\rho(c)>0$ is some constant that approaches $0$ as $c$ approaches $0$. With the choice $F(c) = \max\{2c,\rho(c),\Theta(c)\}$, inequalities \eqref{v1inequalitysum2squares}, \eqref{u1inequalitysum2squares}, \eqref{u0inequalitysum2squares}, \eqref{v0inequalitysum2squares}, and set relation \eqref{Halfdoneu0u1u2v0v1v2proof}, and the discussion stated immediately after \eqref{Halfdoneu0u1u2v0v1v2proof}, the proof follows.
\end{proof}

\end{document}
